%! Skills Focus: Selecting, Implementing, and Communicating Inference Procedures — Unit 7, Lesson 7.10 — Lesson Plan
\documentclass[10pt]{article}
\usepackage{apstats-article}
\usepackage{apstats-boxes}
\usepackage{graphicx}
\graphicspath{{images/}}
\newcommand{\TallMath}[1]{$\displaystyle #1\rule[-1.4em]{0pt}{3.2em}$}
\newcommand{\UnitNumberName}{Unit 7: Inference for Quantitative Data: Means \quad}
\newcommand{\LessonNumberName}{Lesson 7.10: Skills Focus --- Selecting, Implementing, and Communicating Inference Procedures}

\begin{document}
\begin{center}
    {\Large\bfseries \CourseName: \SchoolYear \\
    {\small\bfseries \UnitNumberName \LessonNumberName}}
\end{center}

\begin{tcolorbox}[colback=sky, colframe=navy, boxrule=0.9pt, arc=2mm,
    left=2mm, right=2mm, top=1.5mm, bottom=1.5mm]
    \textbf{Primary Objective:} Students will select the correct inference procedure from all
    six procedures learned in Units 6--7, implement it correctly through all four steps, and
    communicate conclusions in AP-exam language with proper linkage between evidence and claims.
    (Big Ideas: VAR, UNC)
\end{tcolorbox}

\renewcommand{\arraystretch}{1.6}
\begin{skillbox}[Priority Ideas \& Skills]{goldbox}
\begin{minipage}[t]{0.48\linewidth}
    \textbf{AP Skill 1 --- Selecting Statistical Methods}
    \begin{itemize}
        \item Identify an appropriate inference procedure for a given scenario by recognizing
              the parameter type (proportion vs.\ mean) and the number of samples (one vs.\ two)
              (Skills 1.D, 1.E; VAR-6, VAR-7).
        \item Identify when a confidence interval vs.\ a significance test is the
              appropriate tool (Skill 1.D).
        \item State hypotheses and check conditions for the selected procedure
              (Skill 1.F; VAR-6, VAR-7).
    \end{itemize}
    \textbf{AP Skill 3 --- Using Probability and Simulation}
    \begin{itemize}
        \item Interpret a $p$-value in context and connect it to the conclusion
              (Skill 3.D; VAR-6, VAR-7).
        \item Interpret a confidence interval in context (Skill 3.E).
    \end{itemize}
\end{minipage}
\hfill
\begin{minipage}[t]{0.48\linewidth}
    \textbf{AP Skill 4 --- Statistical Argumentation}
    \begin{itemize}
        \item Justify the selection of an inference procedure by citing the parameter and
              sample structure (Skill 4.B).
        \item Verify conditions for the selected procedure and communicate whether they are
              met (Skill 4.C).
        \item Carry out the procedure and report the test statistic and $p$-value (or interval)
              (Skill 4.D).
        \item Provide a conclusion in context using AP-exam language (Skill 4.E).
    \end{itemize}
    \textbf{Key Understandings}
    \begin{itemize}
        \item The two decision axes: (1) proportion or mean? (2) one sample or two?
        \item Every inference procedure follows the same four-step structure.
        \item A CI estimates a parameter; a significance test evaluates a specific claim.
        \item Matched-pairs data $\to$ one-sample $t$-test on the differences.
    \end{itemize}
\end{minipage}
\end{skillbox}

\begin{skillbox}[{Vocabulary, Concepts, \& Theorems}]{greenbox}
\textbf{Procedure Selection Table} --- Ask two questions: (1) Proportion or mean? (2) One sample or two?

\medskip
\renewcommand{\arraystretch}{1.7}
\begin{tabularx}{\linewidth}{@{} p{0.18\linewidth} p{0.16\linewidth} p{0.14\linewidth} X X @{}}
  \rowcolor{sky}
  \textbf{Parameter} & \textbf{\# Samples} & \textbf{Goal} & \textbf{Procedure Name} & \textbf{Statistic / Form} \\
  \hline
  Proportion $p$ & One & CI & One-sample $z$-interval for $p$ & \TallMath{\hat p \pm z^* \sqrt{\tfrac{\hat p(1-\hat p)}{n}}} \\
  Proportion $p$ & One & Test & One-sample $z$-test for $p$ & \TallMath{z = \dfrac{\hat p - p_0}{\sqrt{p_0(1-p_0)/n}}} \\
  $p_1 - p_2$ & Two & CI & Two-sample $z$-interval for $p_1-p_2$ & \TallMath{(\hat p_1 - \hat p_2) \pm z^* SE} \\
  $p_1 - p_2$ & Two & Test & Two-sample $z$-test for $p_1-p_2$ & pooled $\hat p_c$ \\
  Mean $\mu$ & One & CI & One-sample $t$-interval for $\mu$ & \TallMath{\bar x \pm t^* \cdot \tfrac{s}{\sqrt{n}}} \\
  Mean $\mu$ & One & Test & One-sample $t$-test for $\mu$ (or matched pairs) & \TallMath{t = \dfrac{\bar x - \mu_0}{s/\sqrt{n}}} \\
  $\mu_1 - \mu_2$ & Two & CI & Two-sample $t$-interval for $\mu_1-\mu_2$ & \TallMath{(\bar x_1 - \bar x_2) \pm t^* SE} \\
  $\mu_1 - \mu_2$ & Two & Test & Two-sample $t$-test for $\mu_1-\mu_2$ & \TallMath{t = \dfrac{\bar x_1 - \bar x_2}{SE}} \\
\end{tabularx}
\end{skillbox}

\begin{skillbox}[Activate Prior Knowledge \& Spiral Review]{sky}
\begin{minipage}[t]{0.55\linewidth}\vspace{0pt}
    \textbf{Spiral Review (warm-up)}
    \begin{itemize}
        \item All six inference procedures: names, parameter symbols, and which distribution
              ($z$ vs.\ $t$) applies --- Units 6 and 7
        \item Four-step inference structure: State, Plan, Do, Conclude --- Units 6--7
        \item Connecting $p$-values to conclusions; CI interpretation language
    \end{itemize}
\end{minipage}
\hfill
\begin{minipage}[t]{0.38\linewidth}\vspace{0pt}\centering
    % Authored warm-up compiles to target/ — no source PDF to embed here.
\end{minipage}
\end{skillbox}

\begin{skillbox}[Hook]{sky}
\begin{minipage}[t]{0.55\linewidth}
    \textbf{Which Procedure Do I Use? (3--4 min)}\\
    Display this scenario: ``A sports scientist collects resting heart rates for 20 elite
    cyclists and 20 recreational cyclists. She wants to know whether elite cyclists have a
    lower \emph{mean} resting heart rate.''\\[4pt]
    Ask: ``Before we do any math --- what procedure should we use? How do you know?''\\[4pt]
    Facilitate a 60-second Think-Pair-Share. Key reasoning:
    \begin{itemize}
        \item The parameter is a mean, not a proportion.
        \item Two independent groups (not matched pairs).
        \item She wants to test a claim $\to$ significance test.
        \item Conclusion: \textbf{two-sample $t$-test for $\mu_1-\mu_2$}.
    \end{itemize}
\end{minipage}
\hfill
\begin{minipage}[t]{0.40\linewidth}
    \textbf{The Two Decision Questions}
    \begin{enumerate}
        \item \textbf{What is the parameter?}\\
              Proportion ($p$) $\to$ $z$-procedure\\
              Mean ($\mu$) $\to$ $t$-procedure
        \item \textbf{How many samples?}\\
              One group (or matched pairs) $\to$ one-sample\\
              Two independent groups $\to$ two-sample
    \end{enumerate}
    \medskip
    Third question: \textbf{Estimate or test?}\\
    ``What is the range?'' $\to$ CI\\
    ``Is there evidence for a claim?'' $\to$ test
\end{minipage}
\end{skillbox}

\begin{skillbox}[Lesson: The Decision Framework]{sky}
\begin{minipage}[t]{0.48\linewidth}
    \textbf{Step-by-Step Procedure Selection}
    \begin{enumerate}
        \item \textbf{Identify the parameter.}\\
              ``Proportion,'' ``percent,'' ``rate'' $\to$ $p$;\\
              ``Mean,'' ``average'' $\to$ $\mu$.
        \item \textbf{Count the samples.}\\
              One group: one-sample.\\
              Two independent groups: two-sample.\\
              Same subjects measured twice (before/after): matched pairs $\to$ one-sample $t$ on differences.
        \item \textbf{Decide: CI or test?}\\
              Estimate a parameter or find a plausible range $\to$ CI.\\
              Evaluate evidence for a specific claim $\to$ test.
        \item \textbf{Name the procedure} using the exact AP name.
    \end{enumerate}
\end{minipage}
\hfill
\begin{minipage}[t]{0.48\linewidth}
    \textbf{Common Traps}
    \begin{itemize}
        \item ``Before and after'' data: check whether it's \emph{matched pairs} (same subjects)
              or two \emph{independent} samples. Matched pairs $\to$ one-sample $t$ on differences.
        \item Hypotheses use population parameters, not sample statistics:
              $H_0: \mu_1 - \mu_2 = 0$, not $H_0: \bar x_1 - \bar x_2 = 0$.
        \item Two-sample $z$-test: pool ($\hat p_c$) only when testing $H_0: p_1 = p_2$.
              Never pool for a CI.
        \item Two-sample $t$ df: let technology compute Welch df; do not use $n_1+n_2-2$.
    \end{itemize}
\end{minipage}
\end{skillbox}

\begin{skillbox}[Explicit Instruction: AP Four-Step Write-Up for a Significance Test]{sky}
\begin{minipage}[t]{0.46\linewidth}
    \textbf{The Four Steps}
    \begin{enumerate}
        \item \textbf{State.} Name the procedure. State $H_0$ and $H_a$ in terms of the
              population parameter. Define symbols in context.
        \item \textbf{Plan.} List each condition. Verify with numbers from the problem.
              State whether each condition is met.
        \item \textbf{Do.} Show the formula and substitution for the test statistic.
              Report the $p$-value.
        \item \textbf{Conclude.} ``Because $p$-value [$ < $ or $ \ge $] $\alpha = 0.05$,
              I [reject / fail to reject] $H_0$. There [is / is not] sufficient evidence
              to conclude that [restate $H_a$ in context].''
    \end{enumerate}
\end{minipage}
\hfill
\begin{minipage}[t]{0.50\linewidth}
    \textbf{Worked Example: Resting Heart Rate}\\
    \textit{Elite: $n_1=20$, $\bar x_1=52$, $s_1=6$. \;
    Recreational: $n_2=20$, $\bar x_2=61$, $s_2=8$. \\
    Is there evidence elite cyclists have a lower mean HR? ($\alpha=0.05$)}
    \begin{enumerate}
        \item \textbf{State:} Two-sample $t$-test for $\mu_1-\mu_2$; $\mu_1$ = mean HR (elite),
              $\mu_2$ = mean HR (recreational). $H_0: \mu_1-\mu_2=0$; $H_a: \mu_1-\mu_2<0$.
        \item \textbf{Plan:} Random (given). 10\%: $20 \le 10\%$ of all cyclists (assumed large
              population). Normal: $n_1=n_2=20$; no outliers stated; conditions met.
        \item \textbf{Do:} $SE = \sqrt{6^2/20 + 8^2/20} \approx 2.24$;
              $t = (52-61)/2.24 \approx -4.02$; $df \approx 34$ (technology); $p < 0.001$.
        \item \textbf{Conclude:} Because $p < 0.05$, reject $H_0$. There is sufficient evidence
              that the mean HR of elite cyclists is lower than that of recreational cyclists.
    \end{enumerate}
\end{minipage}
\end{skillbox}

\begin{skillbox}[Explicit Instruction: Four-Step Write-Up for a Confidence Interval]{sky}
\begin{minipage}[t]{0.46\linewidth}
    \textbf{The Four Steps for a CI}
    \begin{enumerate}
        \item \textbf{State.} ``I will construct a [level]\% [procedure name].''
              Define the parameter in context.
        \item \textbf{Plan.} Same conditions as the corresponding test. Verify each.
        \item \textbf{Do.} Compute $\text{statistic} \pm (\text{critical value}) \times SE$.
              Report the interval $(L,\, U)$.
        \item \textbf{Conclude.} ``We are [level]\% confident that the [parameter in context]
              is between $L$ and $U$.''
    \end{enumerate}
\end{minipage}
\hfill
\begin{minipage}[t]{0.50\linewidth}
    \textbf{CI vs.\ Significance Test --- When to Use Which}
    \begin{itemize}
        \item Use a \textbf{significance test} when there is a specific null value to evaluate
              (``Is the mean exactly 50?''; ``Are the two proportions equal?'').
        \item Use a \textbf{confidence interval} when the goal is to \emph{estimate} a parameter
              or when the question asks ``how large is the difference?''
        \item A CI can substitute for a two-sided test: if the null value falls outside the CI,
              reject $H_0$.
        \item AP rubric note: when a question says ``use a confidence interval to determine
              whether,'' compute the CI and check whether the null value is captured.
    \end{itemize}
\end{minipage}
\end{skillbox}

\begin{skillbox}[Active Monitoring]{redbox}
\begin{minipage}[t]{0.48\linewidth}
    \textbf{Circulate and Check}
    \begin{itemize}
        \item Students name the correct procedure and can articulate both decision criteria.
        \item Hypotheses use population parameters, not sample statistics.
        \item Conditions are verified with numbers, not just stated as ``met.''
        \item Conclusions connect $p$-value to $\alpha$ explicitly and restate $H_a$ in context.
        \item Students distinguish matched pairs from two independent samples.
    \end{itemize}
\end{minipage}
\hfill
\begin{minipage}[t]{0.48\linewidth}
    \textbf{Cold-Call Prompts}
    \begin{itemize}
        \item ``What two questions do you ask to identify the procedure?''
        \item ``Is this matched pairs or two independent samples? How do you know?''
        \item ``Your $p$-value is 0.08 and $\alpha=0.05$ --- what is your conclusion?
              What if $\alpha=0.10$?''
        \item ``The 95\% CI is $(1.2,\; 4.8)$. Does it suggest a real difference? How?''
    \end{itemize}
\end{minipage}
\end{skillbox}

\begin{skillbox}[Group Work \& Differentiation]{redbox}
\begin{multicols}{3}
    \textbf{Tier R --- Remediate}
    \begin{itemize}
        \item Provide the selection table as a reference.
        \item Activity: identify the procedure and list the conditions only (no calculation).
        \item Focus on the two decision questions.
    \end{itemize}
    \columnbreak
    \textbf{Tier A --- Approaching}
    \begin{itemize}
        \item Complete all three activity scenarios (select, set up, carry out, conclude).
        \item Write full four-step responses in AP-exam language.
    \end{itemize}
    \columnbreak
    \textbf{Tier E --- Extension}
    \begin{itemize}
        \item Complete all scenarios plus critique a provided student AP response.
        \item Identify and correct two errors in procedure choice, conditions, or conclusion.
    \end{itemize}
\end{multicols}
\end{skillbox}

\begin{skillbox}[Individual Work \& Assessment]{redbox}
\begin{minipage}[t]{0.48\linewidth}
    \textbf{Exit Ticket (5 min)}\\
    Two brief scenarios: students name the correct procedure and state one key condition.
    \begin{enumerate}
        \item A researcher surveys 150 randomly selected voters and wants to estimate the
              true proportion who support a ballot measure.
        \item A teacher assigns 32 students to two tutoring methods and tests whether mean
              post-test scores differ between methods.
    \end{enumerate}
\end{minipage}
\hfill
\begin{minipage}[t]{0.48\linewidth}
    \textbf{AP-Style Check}\\
    \textit{Resting pulse rates recorded for 30 randomly selected women and 30 randomly selected
    men. Test whether mean rates differ. Which procedure?}
    \begin{enumerate}[label=(\Alph*)]
        \item One-sample $t$-test for $\mu$
        \item Matched-pairs $t$-test for $\mu_d$
        \item Two-sample $t$-test for $\mu_1-\mu_2$ \textcolor{keyred}{\textbf{$\leftarrow$ correct}}
        \item Two-sample $z$-test for $p_1-p_2$
    \end{enumerate}
\end{minipage}
\end{skillbox}

\begin{skillbox}[Reinforcement \& Extension]{goldbox}
\begin{minipage}[t]{0.48\linewidth}
    \textbf{Homework Overview}
    \begin{itemize}
        \item Five AP-style problems spanning all six inference procedures; at least one
              requires a complete four-step write-up.
        \item One problem asks students to critique a provided response and identify errors.
        \item Extension question: when is a CI more informative than a significance test?
    \end{itemize}
\end{minipage}
\hfill
\begin{minipage}[t]{0.48\linewidth}
    \textbf{Spiral Review}
    \begin{itemize}
        \item All six procedure names, parameters, and distributions (Units 6--7)
        \item Four-step write-up for both CIs and significance tests
        \item $p$-value interpretation and conclusion language
    \end{itemize}
    \textbf{Preview}\\
    Unit 8 introduces chi-square tests --- a new class of procedures for categorical data
    with more than two categories. The same four-step structure applies.
\end{minipage}
\end{skillbox}

\end{document}
